$\mathbf{x}_{j}'=[x_{j1},x_{j2},\cdots,x_{jp}],j=1,2,\cdots,n$ 和
\[
\bar{\mathbf{x}}=\frac{1}{n}\sum_{j=1}^{n}\mathbf{x}_{j}
\]

4.16\quad 设 $\mathbf{X}_{1},\mathbf{X}_{2},\mathbf{X}_{3}$ 和 $\mathbf{X}_{4}$ 是独立 $N_{p}(\boldsymbol{\mu},\boldsymbol{\Sigma})$ 随机向量.

(a) 对以下两个随机向量求边缘分布:
\[
\mathbf{V}_{1}=\frac{1}{4}\mathbf{X}_{1}-\frac{1}{4}\mathbf{X}_{2}+\frac{1}{4}\mathbf{X}_{3}-\frac{1}{4}\mathbf{X}_{4}
\]
\[
\mathbf{V}_{2}=\frac{1}{4}\mathbf{X}_{1}+\frac{1}{4}\mathbf{X}_{2}-\frac{1}{4}\mathbf{X}_{3}-\frac{1}{4}\mathbf{X}_{4}
\]

(b) 求 (a) 中定义的随机向量 $\mathbf{V}_{1}$ 和 $\mathbf{V}_{2}$ 的联合密度.

4.17\quad 设 $\mathbf{X}_{1},\mathbf{X}_{2},\mathbf{X}_{3},\mathbf{X}_{4}$ 和 $\mathbf{X}_{5}$ 是有均值向量 $\boldsymbol{\mu}$ 和协方差矩阵 $\boldsymbol{\Sigma}$ 的独立同分布随机向量, 试对这些随机向量的两个线性组合
\[
\frac{1}{5}\mathbf{X}_{1}+\frac{1}{5}\mathbf{X}_{2}+\frac{1}{5}\mathbf{X}_{3}+\frac{1}{5}\mathbf{X}_{4}+\frac{1}{5}\mathbf{X}_{5}
\]
和
\[
\mathbf{X}_{1}-\mathbf{X}_{2}+\mathbf{X}_{3}-\mathbf{X}_{4}+\mathbf{X}_{5}
\]
分别求均值向量和协方差矩阵, 并求这两个线性组合之间的协方差.

4.18\quad 试根据来自二维正态总体的随机样本
\[
\mathbf{X}=\begin{bmatrix}3&6\\4&4\\5&7\\4&7\end{bmatrix}
\]
求 $2\times 1$ 均值向量 $\boldsymbol{\mu}$ 和 $2\times 2$ 协方差矩阵 $\boldsymbol{\Sigma}$ 的极大似然估计.

4.19\quad 设 $\mathbf{X}_{1},\mathbf{X}_{2},\cdots,\mathbf{X}_{20}$ 为来自总体 $N_{6}(\boldsymbol{\mu},\boldsymbol{\Sigma})$, 容量 $n=20$ 的一个随机样本, 试对下列分布给出完整说明:

(a) $(\mathbf{X}_{1}-\boldsymbol{\mu})'\boldsymbol{\Sigma}^{-1}(\mathbf{X}_{1}-\boldsymbol{\mu})$ 的分布.

(b) $\bar{\mathbf{X}}$ 和 $\sqrt{n}(\bar{\mathbf{X}}-\boldsymbol{\mu})$ 的分布.

(c) $(n-1)\mathbf{S}$ 的分布.

4.20\quad 对练习 4.19 中的随机变量 $\mathbf{X}_{1},\mathbf{X}_{2},\cdots,\mathbf{X}_{20}$, 详细说明以下情况中 $\mathbf{B}(19\mathbf{S})\mathbf{B}'$ 的分布:

(a) $\mathbf{B}=\begin{bmatrix}1&-1/2&-1/2&0&0&0\\0&0&0&-1/2&-1/2&1\end{bmatrix}$

(b) $\mathbf{B}=\begin{bmatrix}1&0&0&0&0&0\\0&0&1&0&0&0\end{bmatrix}$

4.21\quad 设 $\mathbf{X}_{1},\cdots,\mathbf{X}_{60}$ 为来自均值为 $\boldsymbol{\mu}$ 和协方差为 $\boldsymbol{\Sigma}$ 的 4 维正态分布的随机样本, 其容量为 60. 试对下列各项给出完整说明:

(a) $\bar{\mathbf{X}}$ 的分布.

(b) $(\mathbf{X}_{1}-\boldsymbol{\mu})'\boldsymbol{\Sigma}^{-1}(\mathbf{X}_{1}-\boldsymbol{\mu})$ 的分布.

(c) $n(\bar{\mathbf{X}}-\boldsymbol{\mu})'\boldsymbol{\Sigma}^{-1}(\bar{\mathbf{X}}-\boldsymbol{\mu})$ 的分布.

(d) $n(\bar{\mathbf{X}}-\boldsymbol{\mu})'\mathbf{S}^{-1}(\bar{\mathbf{X}}-\boldsymbol{\mu})$ 的近似分布.

4.22\quad 设 $\mathbf{X}_{1},\mathbf{X}_{2},\cdots,\mathbf{X}_{75}$ 为来自均值为 $\boldsymbol{\mu}$, 协方差矩阵为 $\boldsymbol{\Sigma}$ 的总体分布的随机样本, 下面每个统计量的近似分布是什么?
